%% LyX trick_preamble_code_into_believing_that_this_was_created_by_lyx created this file.  For more info, see http://www.lyx.org/.
%% Do not edit unless you really know what you are doing.
\documentclass[a4paper,12pt]{article}
\usepackage[T1]{fontenc}
\usepackage[latin9]{inputenc}

\setlength{\parindent}{3mm}
\usepackage{setspace}
\onehalfspacing

\usepackage{amsmath}
\usepackage{url}
\usepackage{verbatim}
\usepackage[numbers]{natbib}
\usepackage{nomencl}
% the following is useful when we have the old nomencl.sty package
\providecommand{\printnomenclature}{\printglossary}
\providecommand{\makenomenclature}{\makeglossary}
\makenomenclature
\usepackage{varioref}
\usepackage{splitidx}
\makeindex
\newindex[Index]{idx}
\newindex[new]{new}
\newindex{test}
\usepackage{graphicx}
\usepackage{longtable}
\usepackage{xargs}
\usepackage{subscript}
\usepackage{rotating}
\usepackage{listings}
\usepackage{hyperref}
\usepackage{array}
\usepackage{booktabs}
\usepackage{multirow}
\usepackage{hhline}
\usepackage{dcolumn}
\usepackage{pdfpages}
\usepackage{textcomp}
\usepackage{amssymb}
\usepackage{color}
\usepackage[ps,mover]{lyxskak}
\usepackage{tipa}
\usepackage{tipx}
\usepackage{tone}

%%%%%%%%%%%%%%%%%%%%%%%%%%%%%% LyX specific LaTeX commands.
\providecommand{\LyX}{L\kern-.1667em\lower.25em\hbox{Y}\kern-.125emX\@}
\newcommand{\lyxline}[1][1pt]{%
  \par\noindent%
  \rule[.5ex]{\linewidth}{#1}\par}

\newcommand{\lyxarrow}{\leavevmode\,$\triangleright$\,\allowbreak}

%%%%%%%%%%%%%%%%%%%%%%%%%%%%%% User specified LaTeX commands.

% Load refstyle before prettyref so that prettyref wins.
% The real refstyle tests are in test-refstyle-theorems.tex.
\usepackage{refstyle}
\AtBeginDocument{\providecommand\secref[1]{\ref{sec:#1}}}
\usepackage{prettyref}

\def\mycommand{\textquestiondown}

\begin{document}

\title{Title}

\date %stupid stuff
{two days ago}

\maketitle

\tableofcontents
\lstlistoflistings

\noindent This paragraph is not indented.

\section{References\index{References}}

Let's start with simple things: a label~\label{lab:test} and a
reference~\ref{lab:test}; note that I have tested ``unbreakable
space'' without warning (and worse than that, I just tested english
quotes too...).

Of course there are other kind of references, like page
reference~\pageref{lab:test}, but also equation
reference~\eqref{lab:test} (from amsmath package), or varioref's
equivalents~\vref{lab:test} and~\vpageref{lab:test}, without
forgetting pretty references like~\prettyref{lab:test} as well as textual
references like~\nameref{lab:test}.
The command \secref{lab:test} will be parsed in ERT, since LyX
does not support refstyle and prettyref natively at the same time.

\section{Cites\index{Cites}}

Let's start with simple things: a \textbackslash{}cite: \cite[after]{article-crossref}
and two \textbackslash{}cites: \cite[after]{whole-set,article-crossref}

Now the natbib things:

\textbackslash{}citet: \citet[before][after]{article-crossref} \textbackslash{}citet{*}:
\citet*[before][after]{article-crossref} \textbackslash{}Citet: \Citet[before][after]{article-crossref}
\textbackslash{}Citet{*}: \Citet*[before][after]{article-crossref}

\textbackslash{}citep: \citep[before][after]{article-crossref} \textbackslash{}citep{*}:
\citep*[before][after]{article-crossref} \textbackslash{}Citep: \Citep[before][after]{article-crossref}
\textbackslash{}Citep{*}: \Citep*[before][after]{article-crossref}

\textbackslash{}citealt: \citealt[before][after]{article-crossref}
\textbackslash{}citealt{*}: \citealt*[before][after]{article-crossref}
\textbackslash{}Citealt: \Citealt[before][after]{article-crossref}
\textbackslash{}Citealt{*}: \Citealt*[before][after]{article-crossref}

\textbackslash{}citealp: \citealp[before][after]{article-crossref}
\textbackslash{}citealp{*}: \citealp*[before][after]{article-crossref}
\textbackslash{}Citealp: \Citep[before][after]{article-crossref}
\textbackslash{}Citealp{*}: \Citealp*[before][after]{article-crossref}

\textbackslash{}citeauthor: \citeauthor[before][after]{article-crossref}
\textbackslash{}citeauthor{*}: \citeauthor*[before][after]{article-crossref}
\textbackslash{}Citeauthor: \Citeauthor[before][after]{article-crossref}
\textbackslash{}Citeauthor{*}: \Citeauthor*[before][after]{article-crossref}

\textbackslash{}citeyear: \citeyear[before][after]{article-crossref}

\textbackslash{}citeyearpar: \citeyearpar[before][after]{article-crossref} 

\textbackslash{}nocite: \nocite{article-crossref}

% Remove duplicate call of \bibliography since LaTeX throws an error.
%\bibliographystyle{unsrt}
% \bibliography{xampl}

With \textbackslash{}nocite\{{*}\}:
\bibliographystyle{unsrt}
\nocite{*}
\bibliography{xampl}

From bug 7306:
\citet{gur+04} have demonstrated that authors should not have umlauts
in their names. It's OK to use them in the man\"{u}script, though.
\begin{thebibliography}{References}
\bibitem[M\"{o}stl et~al.(2010)]{Mostl2010} G�rkan, M.~A., Freitag,
M., \& Rasio, F.~A. 2004, ApJ, 604, 632
\bibitem[{{G\"{u}rkan et~al.}(2004)}]{gur+04} G�rkan, M.~A., Freitag,
M., \& Rasio, F.~A. 2004, ApJ, 604, 632\end{thebibliography}


\section{Input files\index{Input files}}

We can input files too, like this 
\subsection{\label{sub:External-Subsection}External Subsection}

This is a small dummy child document to show how files can be inserted
to another document. Here are some accented characters to make sure
the encoding is passed to included files: ��
, or with the include
variant (tests quoting as well):
\include{\string"Dummy\space Document\string".tex}
\include{\string"Dummy Document\string".tex}
\include{\string"Dummy\string~Document\string"}
\include{\string"Dummy\string~Document\string".tex}
File names containg double quotes are not legal on windows, so no test file
is provided. Nevertheless, they are legal on unix, so we test at least the
parsing. File names containing spaces can not be handled by automake
(make dist would fail), so no test file is provided as well.
Since the files do not exist, the variants without extension would
not work (tex2lyx can't deduce the correct extension of nonexisting files).
\include{\string"Dummy\string"Document\string".tex}

We can also import chess diagrams:

\loadgame{../../../lib/examples/iecc05}\showboard

Spreadsheets:

\def\inputGnumericTable{}%%%%%%%%%%%%%%%%%%%%%%%%%%%%%%%%%%%%%%%%%%%%%%%%%%%%%%%%%%%%%%%%%%%%%%
%%                                                                  %%
%%  This is the header of a LaTeX2e file exported from Gnumeric.    %%
%%                                                                  %%
%%  This file can be compiled as it stands or included in another   %%
%%  LaTeX document. The table is based on the longtable package so  %%
%%  the longtable options (headers, footers...) can be set in the   %%
%%  preamble section below (see PRAMBLE).                           %%
%%                                                                  %%
%%  To include the file in another, the following two lines must be %%
%%  in the including file:                                          %%
%%        \def\inputGnumericTable{}                                 %%
%%  at the beginning of the file and:                               %%
%%        \input{name-of-this-file.tex}                             %%
%%  where the table is to be placed. Note also that the including   %%
%%  file must use the following packages for the table to be        %%
%%  rendered correctly:                                             %%
%%    \usepackage[latin1]{inputenc}                                 %%
%%    \usepackage{color}                                            %%
%%    \usepackage{array}                                            %%
%%    \usepackage{longtable}                                        %%
%%    \usepackage{calc}                                             %%
%%    \usepackage{multirow}                                         %%
%%    \usepackage{hhline}                                           %%
%%    \usepackage{ifthen}                                           %%
%%  optionally (for landscape tables embedded in another document): %%
%%    \usepackage{lscape}                                           %%
%%                                                                  %%
%%%%%%%%%%%%%%%%%%%%%%%%%%%%%%%%%%%%%%%%%%%%%%%%%%%%%%%%%%%%%%%%%%%%%%



%%  This section checks if we are begin input into another file or  %%
%%  the file will be compiled alone. First use a macro taken from   %%
%%  the TeXbook ex 7.7 (suggestion of Han-Wen Nienhuys).            %%
\def\ifundefined#1{\expandafter\ifx\csname#1\endcsname\relax}


%%  Check for the \def token for inputed files. If it is not        %%
%%  defined, the file will be processed as a standalone and the     %%
%%  preamble will be used.                                          %%
\ifundefined{inputGnumericTable}

%%  We must be able to close or not the document at the end.        %%
	\def\gnumericTableEnd{\end{document}}


%%%%%%%%%%%%%%%%%%%%%%%%%%%%%%%%%%%%%%%%%%%%%%%%%%%%%%%%%%%%%%%%%%%%%%
%%                                                                  %%
%%  This is the PREAMBLE. Change these values to get the right      %%
%%  paper size and other niceties.                                  %%
%%                                                                  %%
%%%%%%%%%%%%%%%%%%%%%%%%%%%%%%%%%%%%%%%%%%%%%%%%%%%%%%%%%%%%%%%%%%%%%%

	\documentclass[12pt%
			  %,landscape%
                    ]{report}
       \usepackage[latin1]{inputenc}
       \usepackage{fullpage}
       \usepackage{color}
       \usepackage{array}
       \usepackage{longtable}
       \usepackage{calc}
       \usepackage{multirow}
       \usepackage{hhline}
       \usepackage{ifthen}

	\begin{document}


%%  End of the preamble for the standalone. The next section is for %%
%%  documents which are included into other LaTeX2e files.          %%
\else

%%  We are not a stand alone document. For a regular table, we will %%
%%  have no preamble and only define the closing to mean nothing.   %%
    \def\gnumericTableEnd{}

%%  If we want landscape mode in an embedded document, comment out  %%
%%  the line above and uncomment the two below. The table will      %%
%%  begin on a new page and run in landscape mode.                  %%
%       \def\gnumericTableEnd{\end{landscape}}
%       \begin{landscape}


%%  End of the else clause for this file being \input.              %%
\fi

%%%%%%%%%%%%%%%%%%%%%%%%%%%%%%%%%%%%%%%%%%%%%%%%%%%%%%%%%%%%%%%%%%%%%%
%%                                                                  %%
%%  The rest is the gnumeric table, except for the closing          %%
%%  statement. Changes below will alter the table's appearance.     %%
%%                                                                  %%
%%%%%%%%%%%%%%%%%%%%%%%%%%%%%%%%%%%%%%%%%%%%%%%%%%%%%%%%%%%%%%%%%%%%%%

\providecommand{\gnumericmathit}[1]{#1} 
%%  Uncomment the next line if you would like your numbers to be in %%
%%  italics if they are italizised in the gnumeric table.           %%
%\renewcommand{\gnumericmathit}[1]{\mathit{#1}}
\providecommand{\gnumericPB}[1]%
{\let\gnumericTemp=\\#1\let\\=\gnumericTemp\hspace{0pt}}
 \ifundefined{gnumericTableWidthDefined}
        \newlength{\gnumericTableWidth}
        \newlength{\gnumericTableWidthComplete}
        \newlength{\gnumericMultiRowLength}
        \global\def\gnumericTableWidthDefined{}
 \fi
%% The following setting protects this code from babel shorthands.  %%
 \ifthenelse{\isundefined{\languageshorthands}}{}{\languageshorthands{english}}
%%  The default table format retains the relative column widths of  %%
%%  gnumeric. They can easily be changed to c, r or l. In that case %%
%%  you may want to comment out the next line and uncomment the one %%
%%  thereafter                                                      %%
\providecommand\gnumbox{\makebox[0pt]}
%%\providecommand\gnumbox[1][]{\makebox}

%% to adjust positions in multirow situations                       %%
\setlength{\bigstrutjot}{\jot}
\setlength{\extrarowheight}{\doublerulesep}

%%  The \setlongtables command keeps column widths the same across  %%
%%  pages. Simply comment out next line for varying column widths.  %%
\setlongtables

\setlength\gnumericTableWidth{%
	100pt+%
	53pt+%
	40pt+%
	53pt+%
0pt}
\def\gumericNumCols{4}
\setlength\gnumericTableWidthComplete{\gnumericTableWidth+%
         \tabcolsep*\gumericNumCols*2+\arrayrulewidth*\gumericNumCols}
\ifthenelse{\lengthtest{\gnumericTableWidthComplete > \linewidth}}%
         {\def\gnumericScale{\ratio{\linewidth-%
                        \tabcolsep*\gumericNumCols*2-%
                        \arrayrulewidth*\gumericNumCols}%
{\gnumericTableWidth}}}%
{\def\gnumericScale{1}}

%%%%%%%%%%%%%%%%%%%%%%%%%%%%%%%%%%%%%%%%%%%%%%%%%%%%%%%%%%%%%%%%%%%%%%
%%                                                                  %%
%% The following are the widths of the various columns. We are      %%
%% defining them here because then they are easier to change.       %%
%% Depending on the cell formats we may use them more than once.    %%
%%                                                                  %%
%%%%%%%%%%%%%%%%%%%%%%%%%%%%%%%%%%%%%%%%%%%%%%%%%%%%%%%%%%%%%%%%%%%%%%

\ifthenelse{\isundefined{\gnumericColB}}{\newlength{\gnumericColB}}{}\settowidth{\gnumericColB}{\begin{tabular}{@{}p{100pt*\gnumericScale}@{}}x\end{tabular}}
\ifthenelse{\isundefined{\gnumericColC}}{\newlength{\gnumericColC}}{}\settowidth{\gnumericColC}{\begin{tabular}{@{}p{53pt*\gnumericScale}@{}}x\end{tabular}}
\ifthenelse{\isundefined{\gnumericColD}}{\newlength{\gnumericColD}}{}\settowidth{\gnumericColD}{\begin{tabular}{@{}p{40pt*\gnumericScale}@{}}x\end{tabular}}
\ifthenelse{\isundefined{\gnumericColE}}{\newlength{\gnumericColE}}{}\settowidth{\gnumericColE}{\begin{tabular}{@{}p{53pt*\gnumericScale}@{}}x\end{tabular}}

\begin{longtable}[c]{%
	b{\gnumericColB}%
	b{\gnumericColC}%
	b{\gnumericColD}%
	b{\gnumericColE}%
	}

%%%%%%%%%%%%%%%%%%%%%%%%%%%%%%%%%%%%%%%%%%%%%%%%%%%%%%%%%%%%%%%%%%%%%%
%%  The longtable options. (Caption, headers... see Goosens, p.124) %%
%	\caption{The Table Caption.}             \\	%
% \hline	% Across the top of the table.
%%  The rest of these options are table rows which are placed on    %%
%%  the first, last or every page. Use \multicolumn if you want.    %%

%%  Header for the first page.                                      %%
%	\multicolumn{4}{c}{The First Header} \\ \hline 
%	\multicolumn{1}{c}{colTag}	%Column 1
%	&\multicolumn{1}{c}{colTag}	%Column 2
%	&\multicolumn{1}{c}{colTag}	%Column 3
%	&\multicolumn{1}{c}{colTag}	\\ \hline %Last column
%	\endfirsthead

%%  The running header definition.                                  %%
%	\hline
%	\multicolumn{4}{l}{\ldots\small\slshape continued} \\ \hline
%	\multicolumn{1}{c}{colTag}	%Column 1
%	&\multicolumn{1}{c}{colTag}	%Column 2
%	&\multicolumn{1}{c}{colTag}	%Column 3
%	&\multicolumn{1}{c}{colTag}	\\ \hline %Last column
%	\endhead

%%  The running footer definition.                                  %%
%	\hline
%	\multicolumn{4}{r}{\small\slshape continued\ldots} \\
%	\endfoot

%%  The ending footer definition.                                   %%
%	\multicolumn{4}{c}{That's all folks} \\ \hline 
%	\endlastfoot
%%%%%%%%%%%%%%%%%%%%%%%%%%%%%%%%%%%%%%%%%%%%%%%%%%%%%%%%%%%%%%%%%%%%%%

\hhline{|----}
	 \multicolumn{4}{|p{	\gnumericColB+%
	\gnumericColC+%
	\gnumericColD+%
	\gnumericColE+%
	\tabcolsep*2*3}|}%
	{\gnumericPB{\centering}\gnumbox{{\color[rgb]{1.00,0.00,0.00} Long spreadsheet for LyX import}}}
\\
	 \multicolumn{4}{|p{	\gnumericColB+%
	\gnumericColC+%
	\gnumericColD+%
	\gnumericColE+%
	\tabcolsep*2*3}|}%
	{\gnumericPB{\centering}\gnumbox{{\color[rgb]{1.00,0.00,0.00} \textit{longsheet.gnumeric}}}}
\\
\hhline{|-|-|-|-|}
	 \multicolumn{1}{|p{\gnumericColB}|}%
	{\gnumericPB{\raggedright}\gnumbox[l]{{\color[rgb]{0.00,0.00,0.50} \textbf{Item}}}}
	&\multicolumn{1}{p{\gnumericColC}|}%
	{\gnumericPB{\raggedleft}\gnumbox[r]{{\color[rgb]{0.00,0.00,0.50} \textbf{Price}}}}
	&\multicolumn{1}{p{\gnumericColD}|}%
	{\gnumericPB{\raggedleft}\gnumbox[r]{{\color[rgb]{0.00,0.00,0.50} \textbf{\#}}}}
	&\multicolumn{1}{p{\gnumericColE}|}%
	{\gnumericPB{\raggedleft}\gnumbox[r]{{\color[rgb]{0.00,0.00,0.50} \textbf{Total}}}}
\\
\hhline{|-|-|-|-|}
	 \multicolumn{1}{|p{\gnumericColB}}%
	{\gnumericPB{\raggedright}\gnumbox[l]{{\color[rgb]{0.00,0.50,0.00} \textit{Stuff}}}}
	&
	&
	&\multicolumn{1}{p{\gnumericColE}|}%
	{}
\\
	 \multicolumn{1}{|p{\gnumericColB}}%
	{\gnumericPB{\raggedright}\gnumbox[l]{scissors}}
	&\gnumericPB{\raggedleft}\gnumbox[r]{8}
	&\gnumericPB{\raggedleft}\gnumbox[r]{30}
	&\multicolumn{1}{p{\gnumericColE}|}%
	{\gnumericPB{\raggedleft}\gnumbox[r]{240}}
\\
	 \multicolumn{1}{|p{\gnumericColB}}%
	{\gnumericPB{\raggedright}\gnumbox[l]{pens}}
	&\gnumericPB{\raggedleft}\gnumbox[r]{1}
	&\gnumericPB{\raggedleft}\gnumbox[r]{30}
	&\multicolumn{1}{p{\gnumericColE}|}%
	{\gnumericPB{\raggedleft}\gnumbox[r]{30}}
\\
	 \multicolumn{1}{|p{\gnumericColB}}%
	{\gnumericPB{\raggedright}\gnumbox[l]{pencils}}
	&\gnumericPB{\raggedleft}\gnumbox[r]{1}
	&\gnumericPB{\raggedleft}\gnumbox[r]{30}
	&\multicolumn{1}{p{\gnumericColE}|}%
	{\gnumericPB{\raggedleft}\gnumbox[r]{30}}
\\
	 \multicolumn{1}{|p{\gnumericColB}}%
	{}
	&
	&
	&\multicolumn{1}{p{\gnumericColE}|}%
	{}
\\
	 \multicolumn{1}{|p{\gnumericColB}}%
	{\gnumericPB{\raggedright}\gnumbox[l]{radio}}
	&\gnumericPB{\raggedleft}\gnumbox[r]{105}
	&\gnumericPB{\raggedleft}\gnumbox[r]{1}
	&\multicolumn{1}{p{\gnumericColE}|}%
	{\gnumericPB{\raggedleft}\gnumbox[r]{105}}
\\
	 \multicolumn{1}{|p{\gnumericColB}}%
	{\gnumericPB{\raggedright}\gnumbox[l]{camera}}
	&\gnumericPB{\raggedleft}\gnumbox[r]{1200}
	&\gnumericPB{\raggedleft}\gnumbox[r]{2}
	&\multicolumn{1}{p{\gnumericColE}|}%
	{\gnumericPB{\raggedleft}\gnumbox[r]{2400}}
\\
	 \multicolumn{1}{|p{\gnumericColB}}%
	{\gnumericPB{\raggedright}\gnumbox[l]{telephoto lens}}
	&\gnumericPB{\raggedleft}\gnumbox[r]{2000}
	&\gnumericPB{\raggedleft}\gnumbox[r]{1}
	&\multicolumn{1}{p{\gnumericColE}|}%
	{\gnumericPB{\raggedleft}\gnumbox[r]{2000}}
\\
	 \multicolumn{1}{|p{\gnumericColB}}%
	{\gnumericPB{\raggedright}\gnumbox[l]{macro lens}}
	&\gnumericPB{\raggedleft}\gnumbox[r]{300}
	&\gnumericPB{\raggedleft}\gnumbox[r]{2}
	&\multicolumn{1}{p{\gnumericColE}|}%
	{\gnumericPB{\raggedleft}\gnumbox[r]{600}}
\\
	 \multicolumn{1}{|p{\gnumericColB}}%
	{\gnumericPB{\raggedright}\gnumbox[l]{flash}}
	&\gnumericPB{\raggedleft}\gnumbox[r]{150}
	&\gnumericPB{\raggedleft}\gnumbox[r]{4}
	&\multicolumn{1}{p{\gnumericColE}|}%
	{\gnumericPB{\raggedleft}\gnumbox[r]{600}}
\\
	 \multicolumn{1}{|p{\gnumericColB}}%
	{\gnumericPB{\raggedright}\gnumbox[l]{memory card 8G}}
	&\gnumericPB{\raggedleft}\gnumbox[r]{40}
	&\gnumericPB{\raggedleft}\gnumbox[r]{16}
	&\multicolumn{1}{p{\gnumericColE}|}%
	{\gnumericPB{\raggedleft}\gnumbox[r]{640}}
\\
	 \multicolumn{1}{|p{\gnumericColB}}%
	{}
	&
	&
	&\multicolumn{1}{p{\gnumericColE}|}%
	{}
\\
	 \multicolumn{1}{|p{\gnumericColB}}%
	{\gnumericPB{\raggedright}\gnumbox[l]{computer}}
	&\gnumericPB{\raggedleft}\gnumbox[r]{300}
	&\gnumericPB{\raggedleft}\gnumbox[r]{30}
	&\multicolumn{1}{p{\gnumericColE}|}%
	{\gnumericPB{\raggedleft}\gnumbox[r]{9000}}
\\
	 \multicolumn{1}{|p{\gnumericColB}}%
	{\gnumericPB{\raggedright}\gnumbox[l]{os}}
	&\gnumericPB{\raggedleft}\gnumbox[r]{0}
	&\gnumericPB{\raggedleft}\gnumbox[r]{30}
	&\multicolumn{1}{p{\gnumericColE}|}%
	{\gnumericPB{\raggedleft}\gnumbox[r]{0}}
\\
	 \multicolumn{1}{|p{\gnumericColB}}%
	{\gnumericPB{\raggedright}\gnumbox[l]{sw packages}}
	&\gnumericPB{\raggedleft}\gnumbox[r]{0}
	&\gnumericPB{\raggedleft}\gnumbox[r]{2600}
	&\multicolumn{1}{p{\gnumericColE}|}%
	{\gnumericPB{\raggedleft}\gnumbox[r]{0}}
\\
	 \multicolumn{1}{|p{\gnumericColB}}%
	{\gnumericPB{\raggedright}\gnumbox[l]{sw docs}}
	&\gnumericPB{\raggedleft}\gnumbox[r]{0}
	&\gnumericPB{\raggedleft}\gnumbox[r]{2600}
	&\multicolumn{1}{p{\gnumericColE}|}%
	{\gnumericPB{\raggedleft}\gnumbox[r]{0}}
\\
	 \multicolumn{1}{|p{\gnumericColB}}%
	{\gnumericPB{\raggedright}\gnumbox[l]{license mgmt}}
	&\gnumericPB{\raggedleft}\gnumbox[r]{0}
	&\gnumericPB{\raggedleft}\gnumbox[r]{0}
	&\multicolumn{1}{p{\gnumericColE}|}%
	{\gnumericPB{\raggedleft}\gnumbox[r]{0}}
\\
	 \multicolumn{1}{|p{\gnumericColB}}%
	{\gnumericPB{\raggedright}\gnumbox[l]{sw upgrades}}
	&\gnumericPB{\raggedleft}\gnumbox[r]{0}
	&\gnumericPB{\raggedleft}\gnumbox[r]{2600}
	&\multicolumn{1}{p{\gnumericColE}|}%
	{\gnumericPB{\raggedleft}\gnumbox[r]{0}}
\\
	 \multicolumn{1}{|p{\gnumericColB}}%
	{}
	&
	&
	&\multicolumn{1}{p{\gnumericColE}|}%
	{}
\\
	 \multicolumn{1}{|p{\gnumericColB}}%
	{\gnumericPB{\raggedright}\gnumbox[l]{{\color[rgb]{0.00,0.50,0.00} \textit{Clothes}}}}
	&
	&
	&\multicolumn{1}{p{\gnumericColE}|}%
	{}
\\
	 \multicolumn{1}{|p{\gnumericColB}}%
	{\gnumericPB{\raggedright}\gnumbox[l]{socks}}
	&\gnumericPB{\raggedleft}\gnumbox[r]{1}
	&\gnumericPB{\raggedleft}\gnumbox[r]{14}
	&\multicolumn{1}{p{\gnumericColE}|}%
	{\gnumericPB{\raggedleft}\gnumbox[r]{14}}
\\
	 \multicolumn{1}{|p{\gnumericColB}}%
	{\gnumericPB{\raggedright}\gnumbox[l]{hats}}
	&\gnumericPB{\raggedleft}\gnumbox[r]{8}
	&\gnumericPB{\raggedleft}\gnumbox[r]{3}
	&\multicolumn{1}{p{\gnumericColE}|}%
	{\gnumericPB{\raggedleft}\gnumbox[r]{24}}
\\
	 \multicolumn{1}{|p{\gnumericColB}}%
	{\gnumericPB{\raggedright}\gnumbox[l]{jeans}}
	&\gnumericPB{\raggedleft}\gnumbox[r]{30}
	&\gnumericPB{\raggedleft}\gnumbox[r]{2}
	&\multicolumn{1}{p{\gnumericColE}|}%
	{\gnumericPB{\raggedleft}\gnumbox[r]{60}}
\\
	 \multicolumn{1}{|p{\gnumericColB}}%
	{\gnumericPB{\raggedright}\gnumbox[l]{black shoes}}
	&\gnumericPB{\raggedleft}\gnumbox[r]{45}
	&\gnumericPB{\raggedleft}\gnumbox[r]{2}
	&\multicolumn{1}{p{\gnumericColE}|}%
	{\gnumericPB{\raggedleft}\gnumbox[r]{90}}
\\
	 \multicolumn{1}{|p{\gnumericColB}}%
	{\gnumericPB{\raggedright}\gnumbox[l]{brown shoes}}
	&\gnumericPB{\raggedleft}\gnumbox[r]{42}
	&\gnumericPB{\raggedleft}\gnumbox[r]{6}
	&\multicolumn{1}{p{\gnumericColE}|}%
	{\gnumericPB{\raggedleft}\gnumbox[r]{252}}
\\
	 \multicolumn{1}{|p{\gnumericColB}}%
	{\gnumericPB{\raggedright}\gnumbox[l]{jacket}}
	&\gnumericPB{\raggedleft}\gnumbox[r]{50}
	&\gnumericPB{\raggedleft}\gnumbox[r]{1}
	&\multicolumn{1}{p{\gnumericColE}|}%
	{\gnumericPB{\raggedleft}\gnumbox[r]{50}}
\\
	 \multicolumn{1}{|p{\gnumericColB}}%
	{\gnumericPB{\raggedright}\gnumbox[l]{shirts}}
	&\gnumericPB{\raggedleft}\gnumbox[r]{22}
	&\gnumericPB{\raggedleft}\gnumbox[r]{14}
	&\multicolumn{1}{p{\gnumericColE}|}%
	{\gnumericPB{\raggedleft}\gnumbox[r]{308}}
\\
	 \multicolumn{1}{|p{\gnumericColB}}%
	{\gnumericPB{\raggedright}\gnumbox[l]{t-shirts}}
	&\gnumericPB{\raggedleft}\gnumbox[r]{5}
	&\gnumericPB{\raggedleft}\gnumbox[r]{6}
	&\multicolumn{1}{p{\gnumericColE}|}%
	{\gnumericPB{\raggedleft}\gnumbox[r]{30}}
\\
	 \multicolumn{1}{|p{\gnumericColB}}%
	{\gnumericPB{\raggedright}\gnumbox[l]{funny t-shirts}}
	&\gnumericPB{\raggedleft}\gnumbox[r]{8}
	&\gnumericPB{\raggedleft}\gnumbox[r]{4}
	&\multicolumn{1}{p{\gnumericColE}|}%
	{\gnumericPB{\raggedleft}\gnumbox[r]{32}}
\\
	 \multicolumn{1}{|p{\gnumericColB}}%
	{}
	&
	&
	&\multicolumn{1}{p{\gnumericColE}|}%
	{}
\\
	 \multicolumn{1}{|p{\gnumericColB}}%
	{}
	&
	&
	&\multicolumn{1}{p{\gnumericColE}|}%
	{}
\\
	 \multicolumn{1}{|p{\gnumericColB}}%
	{\gnumericPB{\raggedright}\gnumbox[l]{{\color[rgb]{0.00,0.50,0.00} \textit{Food}}}}
	&
	&
	&\multicolumn{1}{p{\gnumericColE}|}%
	{}
\\
	 \multicolumn{1}{|p{\gnumericColB}}%
	{\gnumericPB{\raggedright}\gnumbox[l]{cheese}}
	&\gnumericPB{\raggedleft}\gnumbox[r]{3}
	&\gnumericPB{\raggedleft}\gnumbox[r]{3}
	&\multicolumn{1}{p{\gnumericColE}|}%
	{\gnumericPB{\raggedleft}\gnumbox[r]{9}}
\\
	 \multicolumn{1}{|p{\gnumericColB}}%
	{\gnumericPB{\raggedright}\gnumbox[l]{beer}}
	&\gnumericPB{\raggedleft}\gnumbox[r]{1}
	&\gnumericPB{\raggedleft}\gnumbox[r]{48}
	&\multicolumn{1}{p{\gnumericColE}|}%
	{\gnumericPB{\raggedleft}\gnumbox[r]{48}}
\\
	 \multicolumn{1}{|p{\gnumericColB}}%
	{\gnumericPB{\raggedright}\gnumbox[l]{bread}}
	&\gnumericPB{\raggedleft}\gnumbox[r]{1}
	&\gnumericPB{\raggedleft}\gnumbox[r]{3}
	&\multicolumn{1}{p{\gnumericColE}|}%
	{\gnumericPB{\raggedleft}\gnumbox[r]{3}}
\\
	 \multicolumn{1}{|p{\gnumericColB}}%
	{\gnumericPB{\raggedright}\gnumbox[l]{butter}}
	&\gnumericPB{\raggedleft}\gnumbox[r]{3}
	&\gnumericPB{\raggedleft}\gnumbox[r]{1}
	&\multicolumn{1}{p{\gnumericColE}|}%
	{\gnumericPB{\raggedleft}\gnumbox[r]{3}}
\\
	 \multicolumn{1}{|p{\gnumericColB}}%
	{\gnumericPB{\raggedright}\gnumbox[l]{garlic}}
	&\gnumericPB{\raggedleft}\gnumbox[r]{2}
	&\gnumericPB{\raggedleft}\gnumbox[r]{5}
	&\multicolumn{1}{p{\gnumericColE}|}%
	{\gnumericPB{\raggedleft}\gnumbox[r]{10}}
\\
	 \multicolumn{1}{|p{\gnumericColB}}%
	{\gnumericPB{\raggedright}\gnumbox[l]{meatballs}}
	&\gnumericPB{\raggedleft}\gnumbox[r]{1}
	&\gnumericPB{\raggedleft}\gnumbox[r]{50}
	&\multicolumn{1}{p{\gnumericColE}|}%
	{\gnumericPB{\raggedleft}\gnumbox[r]{50}}
\\
	 \multicolumn{1}{|p{\gnumericColB}}%
	{\gnumericPB{\raggedright}\gnumbox[l]{mustard}}
	&\gnumericPB{\raggedleft}\gnumbox[r]{13}
	&\gnumericPB{\raggedleft}\gnumbox[r]{1}
	&\multicolumn{1}{p{\gnumericColE}|}%
	{\gnumericPB{\raggedleft}\gnumbox[r]{13}}
\\
	 \multicolumn{1}{|p{\gnumericColB}}%
	{\gnumericPB{\raggedright}\gnumbox[l]{ketchup}}
	&\gnumericPB{\raggedleft}\gnumbox[r]{11}
	&\gnumericPB{\raggedleft}\gnumbox[r]{1}
	&\multicolumn{1}{p{\gnumericColE}|}%
	{\gnumericPB{\raggedleft}\gnumbox[r]{11}}
\\
	 \multicolumn{1}{|p{\gnumericColB}}%
	{\gnumericPB{\raggedright}\gnumbox[l]{salt}}
	&\gnumericPB{\raggedleft}\gnumbox[r]{2}
	&\gnumericPB{\raggedleft}\gnumbox[r]{6}
	&\multicolumn{1}{p{\gnumericColE}|}%
	{\gnumericPB{\raggedleft}\gnumbox[r]{12}}
\\
	 \multicolumn{1}{|p{\gnumericColB}}%
	{\gnumericPB{\raggedright}\gnumbox[l]{pepper}}
	&\gnumericPB{\raggedleft}\gnumbox[r]{3}
	&\gnumericPB{\raggedleft}\gnumbox[r]{2}
	&\multicolumn{1}{p{\gnumericColE}|}%
	{\gnumericPB{\raggedleft}\gnumbox[r]{6}}
\\
	 \multicolumn{1}{|p{\gnumericColB}}%
	{}
	&
	&
	&\multicolumn{1}{p{\gnumericColE}|}%
	{}
\\
	 \multicolumn{1}{|p{\gnumericColB}}%
	{\gnumericPB{\raggedright}\gnumbox[l]{{\color[rgb]{0.00,0.50,0.00} \textit{Animals}}}}
	&
	&
	&\multicolumn{1}{p{\gnumericColE}|}%
	{}
\\
	 \multicolumn{1}{|p{\gnumericColB}}%
	{\gnumericPB{\raggedright}\gnumbox[l]{cows}}
	&\gnumericPB{\raggedleft}\gnumbox[r]{500}
	&\gnumericPB{\raggedleft}\gnumbox[r]{30}
	&\multicolumn{1}{p{\gnumericColE}|}%
	{\gnumericPB{\raggedleft}\gnumbox[r]{15000}}
\\
	 \multicolumn{1}{|p{\gnumericColB}}%
	{\gnumericPB{\raggedright}\gnumbox[l]{sheep}}
	&\gnumericPB{\raggedleft}\gnumbox[r]{300}
	&\gnumericPB{\raggedleft}\gnumbox[r]{800}
	&\multicolumn{1}{p{\gnumericColE}|}%
	{\gnumericPB{\raggedleft}\gnumbox[r]{240000}}
\\
	 \multicolumn{1}{|p{\gnumericColB}}%
	{\gnumericPB{\raggedright}\gnumbox[l]{dogs}}
	&\gnumericPB{\raggedleft}\gnumbox[r]{250}
	&\gnumericPB{\raggedleft}\gnumbox[r]{8}
	&\multicolumn{1}{p{\gnumericColE}|}%
	{\gnumericPB{\raggedleft}\gnumbox[r]{2000}}
\\
	 \multicolumn{1}{|p{\gnumericColB}}%
	{\gnumericPB{\raggedright}\gnumbox[l]{cats}}
	&\gnumericPB{\raggedleft}\gnumbox[r]{5}
	&\gnumericPB{\raggedleft}\gnumbox[r]{3}
	&\multicolumn{1}{p{\gnumericColE}|}%
	{\gnumericPB{\raggedleft}\gnumbox[r]{15}}
\\
	 \multicolumn{1}{|p{\gnumericColB}}%
	{\gnumericPB{\raggedright}\gnumbox[l]{goats}}
	&\gnumericPB{\raggedleft}\gnumbox[r]{350}
	&\gnumericPB{\raggedleft}\gnumbox[r]{60}
	&\multicolumn{1}{p{\gnumericColE}|}%
	{\gnumericPB{\raggedleft}\gnumbox[r]{21000}}
\\
	 \multicolumn{1}{|p{\gnumericColB}}%
	{\gnumericPB{\raggedright}\gnumbox[l]{horse}}
	&\gnumericPB{\raggedleft}\gnumbox[r]{660}
	&\gnumericPB{\raggedleft}\gnumbox[r]{1}
	&\multicolumn{1}{p{\gnumericColE}|}%
	{\gnumericPB{\raggedleft}\gnumbox[r]{660}}
\\
	 \multicolumn{1}{|p{\gnumericColB}}%
	{\gnumericPB{\raggedright}\gnumbox[l]{pigs}}
	&\gnumericPB{\raggedleft}\gnumbox[r]{400}
	&\gnumericPB{\raggedleft}\gnumbox[r]{200}
	&\multicolumn{1}{p{\gnumericColE}|}%
	{\gnumericPB{\raggedleft}\gnumbox[r]{80000}}
\\
	 \multicolumn{1}{|p{\gnumericColB}}%
	{\gnumericPB{\raggedright}\gnumbox[l]{bears}}
	&\gnumericPB{\raggedleft}\gnumbox[r]{1800}
	&\gnumericPB{\raggedleft}\gnumbox[r]{4}
	&\multicolumn{1}{p{\gnumericColE}|}%
	{\gnumericPB{\raggedleft}\gnumbox[r]{7200}}
\\
	 \multicolumn{1}{|p{\gnumericColB}}%
	{\gnumericPB{\raggedright}\gnumbox[l]{bears (polar)}}
	&\gnumericPB{\raggedleft}\gnumbox[r]{5000}
	&\gnumericPB{\raggedleft}\gnumbox[r]{2}
	&\multicolumn{1}{p{\gnumericColE}|}%
	{\gnumericPB{\raggedleft}\gnumbox[r]{10000}}
\\
	 \multicolumn{1}{|p{\gnumericColB}}%
	{\gnumericPB{\raggedright}\gnumbox[l]{giraffes}}
	&\gnumericPB{\raggedleft}\gnumbox[r]{4500}
	&\gnumericPB{\raggedleft}\gnumbox[r]{2}
	&\multicolumn{1}{p{\gnumericColE}|}%
	{\gnumericPB{\raggedleft}\gnumbox[r]{9000}}
\\
	 \multicolumn{1}{|p{\gnumericColB}}%
	{\gnumericPB{\raggedright}\gnumbox[l]{hippos}}
	&\gnumericPB{\raggedleft}\gnumbox[r]{12000}
	&\gnumericPB{\raggedleft}\gnumbox[r]{16}
	&\multicolumn{1}{p{\gnumericColE}|}%
	{\gnumericPB{\raggedleft}\gnumbox[r]{192000}}
\\
	 \multicolumn{1}{|p{\gnumericColB}}%
	{\gnumericPB{\raggedright}\gnumbox[l]{camels}}
	&\gnumericPB{\raggedleft}\gnumbox[r]{800}
	&\gnumericPB{\raggedleft}\gnumbox[r]{12}
	&\multicolumn{1}{p{\gnumericColE}|}%
	{\gnumericPB{\raggedleft}\gnumbox[r]{9600}}
\\
	 \multicolumn{1}{|p{\gnumericColB}}%
	{\gnumericPB{\raggedright}\gnumbox[l]{squirrels}}
	&\gnumericPB{\raggedleft}\gnumbox[r]{8}
	&\gnumericPB{\raggedleft}\gnumbox[r]{16}
	&\multicolumn{1}{p{\gnumericColE}|}%
	{\gnumericPB{\raggedleft}\gnumbox[r]{128}}
\\
	 \multicolumn{1}{|p{\gnumericColB}}%
	{\gnumericPB{\raggedright}\gnumbox[l]{mice}}
	&\gnumericPB{\raggedleft}\gnumbox[r]{1}
	&\gnumericPB{\raggedleft}\gnumbox[r]{45}
	&\multicolumn{1}{p{\gnumericColE}|}%
	{\gnumericPB{\raggedleft}\gnumbox[r]{45}}
\\
	 \multicolumn{1}{|p{\gnumericColB}}%
	{\gnumericPB{\raggedright}\gnumbox[l]{mice (white)}}
	&\gnumericPB{\raggedleft}\gnumbox[r]{2}
	&\gnumericPB{\raggedleft}\gnumbox[r]{133}
	&\multicolumn{1}{p{\gnumericColE}|}%
	{\gnumericPB{\raggedleft}\gnumbox[r]{266}}
\\
	 \multicolumn{1}{|p{\gnumericColB}}%
	{\gnumericPB{\raggedright}\gnumbox[l]{rats}}
	&\gnumericPB{\raggedleft}\gnumbox[r]{2}
	&\gnumericPB{\raggedleft}\gnumbox[r]{20}
	&\multicolumn{1}{p{\gnumericColE}|}%
	{\gnumericPB{\raggedleft}\gnumbox[r]{40}}
\\
	 \multicolumn{1}{|p{\gnumericColB}}%
	{\gnumericPB{\raggedright}\gnumbox[l]{elephants}}
	&\gnumericPB{\raggedleft}\gnumbox[r]{10000}
	&\gnumericPB{\raggedleft}\gnumbox[r]{20}
	&\multicolumn{1}{p{\gnumericColE}|}%
	{\gnumericPB{\raggedleft}\gnumbox[r]{200000}}
\\
	 \multicolumn{1}{|p{\gnumericColB}}%
	{\gnumericPB{\raggedright}\gnumbox[l]{crocodile (s)}}
	&\gnumericPB{\raggedleft}\gnumbox[r]{500}
	&\gnumericPB{\raggedleft}\gnumbox[r]{18}
	&\multicolumn{1}{p{\gnumericColE}|}%
	{\gnumericPB{\raggedleft}\gnumbox[r]{9000}}
\\
	 \multicolumn{1}{|p{\gnumericColB}}%
	{\gnumericPB{\raggedright}\gnumbox[l]{crocodile (m)}}
	&\gnumericPB{\raggedleft}\gnumbox[r]{1500}
	&\gnumericPB{\raggedleft}\gnumbox[r]{8}
	&\multicolumn{1}{p{\gnumericColE}|}%
	{\gnumericPB{\raggedleft}\gnumbox[r]{12000}}
\\
	 \multicolumn{1}{|p{\gnumericColB}}%
	{\gnumericPB{\raggedright}\gnumbox[l]{crocodile (l)}}
	&\gnumericPB{\raggedleft}\gnumbox[r]{3500}
	&\gnumericPB{\raggedleft}\gnumbox[r]{2}
	&\multicolumn{1}{p{\gnumericColE}|}%
	{\gnumericPB{\raggedleft}\gnumbox[r]{7000}}
\\
	 \multicolumn{1}{|p{\gnumericColB}}%
	{\gnumericPB{\raggedright}\gnumbox[l]{crocodile (xl)}}
	&\gnumericPB{\raggedleft}\gnumbox[r]{8000}
	&\gnumericPB{\raggedleft}\gnumbox[r]{1}
	&\multicolumn{1}{p{\gnumericColE}|}%
	{\gnumericPB{\raggedleft}\gnumbox[r]{8000}}
\\
	 \multicolumn{1}{|p{\gnumericColB}}%
	{\gnumericPB{\raggedright}\gnumbox[l]{gnu}}
	&\gnumericPB{\raggedleft}\gnumbox[r]{655}
	&\gnumericPB{\raggedleft}\gnumbox[r]{45}
	&\multicolumn{1}{p{\gnumericColE}|}%
	{\gnumericPB{\raggedleft}\gnumbox[r]{29475}}
\\
	 \multicolumn{1}{|p{\gnumericColB}}%
	{\gnumericPB{\raggedright}\gnumbox[l]{leopards}}
	&\gnumericPB{\raggedleft}\gnumbox[r]{9600}
	&\gnumericPB{\raggedleft}\gnumbox[r]{4}
	&\multicolumn{1}{p{\gnumericColE}|}%
	{\gnumericPB{\raggedleft}\gnumbox[r]{38400}}
\\
	 \multicolumn{1}{|p{\gnumericColB}}%
	{\gnumericPB{\raggedright}\gnumbox[l]{weasel}}
	&\gnumericPB{\raggedleft}\gnumbox[r]{45}
	&\gnumericPB{\raggedleft}\gnumbox[r]{12}
	&\multicolumn{1}{p{\gnumericColE}|}%
	{\gnumericPB{\raggedleft}\gnumbox[r]{540}}
\\
	 \multicolumn{1}{|p{\gnumericColB}}%
	{\gnumericPB{\raggedright}\gnumbox[l]{zebras}}
	&\gnumericPB{\raggedleft}\gnumbox[r]{400}
	&\gnumericPB{\raggedleft}\gnumbox[r]{12}
	&\multicolumn{1}{p{\gnumericColE}|}%
	{\gnumericPB{\raggedleft}\gnumbox[r]{4800}}
\\
	 \multicolumn{1}{|p{\gnumericColB}}%
	{\gnumericPB{\raggedright}\gnumbox[l]{ant eater}}
	&\gnumericPB{\raggedleft}\gnumbox[r]{83}
	&\gnumericPB{\raggedleft}\gnumbox[r]{1}
	&\multicolumn{1}{p{\gnumericColE}|}%
	{\gnumericPB{\raggedleft}\gnumbox[r]{83}}
\\
	 \multicolumn{1}{|p{\gnumericColB}}%
	{\gnumericPB{\raggedright}\gnumbox[l]{ants}}
	&\gnumericPB{\raggedleft}\gnumbox[r]{1}
	&\gnumericPB{\raggedleft}\gnumbox[r]{1300}
	&\multicolumn{1}{p{\gnumericColE}|}%
	{\gnumericPB{\raggedleft}\gnumbox[r]{1300}}
\\
	 \multicolumn{1}{|p{\gnumericColB}}%
	{\gnumericPB{\raggedright}\gnumbox[l]{gorillas}}
	&\gnumericPB{\raggedleft}\gnumbox[r]{12000}
	&\gnumericPB{\raggedleft}\gnumbox[r]{5}
	&\multicolumn{1}{p{\gnumericColE}|}%
	{\gnumericPB{\raggedleft}\gnumbox[r]{60000}}
\\
	 \multicolumn{1}{|p{\gnumericColB}}%
	{\gnumericPB{\raggedright}\gnumbox[l]{monkeys}}
	&\gnumericPB{\raggedleft}\gnumbox[r]{300}
	&\gnumericPB{\raggedleft}\gnumbox[r]{45}
	&\multicolumn{1}{p{\gnumericColE}|}%
	{\gnumericPB{\raggedleft}\gnumbox[r]{13500}}
\\
	 \multicolumn{1}{|p{\gnumericColB}}%
	{\gnumericPB{\raggedright}\gnumbox[l]{lions}}
	&\gnumericPB{\raggedleft}\gnumbox[r]{8000}
	&\gnumericPB{\raggedleft}\gnumbox[r]{6}
	&\multicolumn{1}{p{\gnumericColE}|}%
	{\gnumericPB{\raggedleft}\gnumbox[r]{48000}}
\\
	 \multicolumn{1}{|p{\gnumericColB}}%
	{}
	&
	&
	&\multicolumn{1}{p{\gnumericColE}|}%
	{}
\\
	 \multicolumn{1}{|p{\gnumericColB}}%
	{}
	&
	&
	&\multicolumn{1}{p{\gnumericColE}|}%
	{}
\\
	 \multicolumn{1}{|p{\gnumericColB}}%
	{\gnumericPB{\raggedright}\gnumbox[l]{{\color[rgb]{0.00,0.50,0.00} \textit{Furniture}}}}
	&
	&
	&\multicolumn{1}{p{\gnumericColE}|}%
	{}
\\
	 \multicolumn{1}{|p{\gnumericColB}}%
	{\gnumericPB{\raggedright}\gnumbox[l]{chairs}}
	&\gnumericPB{\raggedleft}\gnumbox[r]{12}
	&\gnumericPB{\raggedleft}\gnumbox[r]{64}
	&\multicolumn{1}{p{\gnumericColE}|}%
	{\gnumericPB{\raggedleft}\gnumbox[r]{768}}
\\
	 \multicolumn{1}{|p{\gnumericColB}}%
	{\gnumericPB{\raggedright}\gnumbox[l]{tables}}
	&\gnumericPB{\raggedleft}\gnumbox[r]{20}
	&\gnumericPB{\raggedleft}\gnumbox[r]{6}
	&\multicolumn{1}{p{\gnumericColE}|}%
	{\gnumericPB{\raggedleft}\gnumbox[r]{120}}
\\
	 \multicolumn{1}{|p{\gnumericColB}}%
	{\gnumericPB{\raggedright}\gnumbox[l]{table (round)}}
	&\gnumericPB{\raggedleft}\gnumbox[r]{33}
	&\gnumericPB{\raggedleft}\gnumbox[r]{1}
	&\multicolumn{1}{p{\gnumericColE}|}%
	{\gnumericPB{\raggedleft}\gnumbox[r]{33}}
\\
	 \multicolumn{1}{|p{\gnumericColB}}%
	{\gnumericPB{\raggedright}\gnumbox[l]{chairs (office)}}
	&\gnumericPB{\raggedleft}\gnumbox[r]{16}
	&\gnumericPB{\raggedleft}\gnumbox[r]{34}
	&\multicolumn{1}{p{\gnumericColE}|}%
	{\gnumericPB{\raggedleft}\gnumbox[r]{544}}
\\
	 \multicolumn{1}{|p{\gnumericColB}}%
	{\gnumericPB{\raggedright}\gnumbox[l]{desks}}
	&\gnumericPB{\raggedleft}\gnumbox[r]{42}
	&\gnumericPB{\raggedleft}\gnumbox[r]{30}
	&\multicolumn{1}{p{\gnumericColE}|}%
	{\gnumericPB{\raggedleft}\gnumbox[r]{1260}}
\\
	 \multicolumn{1}{|p{\gnumericColB}}%
	{\gnumericPB{\raggedright}\gnumbox[l]{big desk}}
	&\gnumericPB{\raggedleft}\gnumbox[r]{150}
	&\gnumericPB{\raggedleft}\gnumbox[r]{1}
	&\multicolumn{1}{p{\gnumericColE}|}%
	{\gnumericPB{\raggedleft}\gnumbox[r]{150}}
\\
	 \multicolumn{1}{|p{\gnumericColB}}%
	{\gnumericPB{\raggedright}\gnumbox[l]{boss chair}}
	&\gnumericPB{\raggedleft}\gnumbox[r]{120}
	&\gnumericPB{\raggedleft}\gnumbox[r]{1}
	&\multicolumn{1}{p{\gnumericColE}|}%
	{\gnumericPB{\raggedleft}\gnumbox[r]{120}}
\\
	 \multicolumn{1}{|p{\gnumericColB}}%
	{\gnumericPB{\raggedright}\gnumbox[l]{desk lamps}}
	&\gnumericPB{\raggedleft}\gnumbox[r]{8}
	&\gnumericPB{\raggedleft}\gnumbox[r]{30}
	&\multicolumn{1}{p{\gnumericColE}|}%
	{\gnumericPB{\raggedleft}\gnumbox[r]{240}}
\\
	 \multicolumn{1}{|p{\gnumericColB}}%
	{\gnumericPB{\raggedright}\gnumbox[l]{wall lamps}}
	&\gnumericPB{\raggedleft}\gnumbox[r]{4}
	&\gnumericPB{\raggedleft}\gnumbox[r]{22}
	&\multicolumn{1}{p{\gnumericColE}|}%
	{\gnumericPB{\raggedleft}\gnumbox[r]{88}}
\\
	 \multicolumn{1}{|p{\gnumericColB}}%
	{\gnumericPB{\raggedright}\gnumbox[l]{standing lamps}}
	&\gnumericPB{\raggedleft}\gnumbox[r]{11}
	&\gnumericPB{\raggedleft}\gnumbox[r]{12}
	&\multicolumn{1}{p{\gnumericColE}|}%
	{\gnumericPB{\raggedleft}\gnumbox[r]{132}}
\\
	 \multicolumn{1}{|p{\gnumericColB}}%
	{\gnumericPB{\raggedright}\gnumbox[l]{water coolers}}
	&\gnumericPB{\raggedleft}\gnumbox[r]{57}
	&\gnumericPB{\raggedleft}\gnumbox[r]{5}
	&\multicolumn{1}{p{\gnumericColE}|}%
	{\gnumericPB{\raggedleft}\gnumbox[r]{285}}
\\
	 \multicolumn{1}{|p{\gnumericColB}}%
	{\gnumericPB{\raggedright}\gnumbox[l]{safe}}
	&\gnumericPB{\raggedleft}\gnumbox[r]{6500}
	&\gnumericPB{\raggedleft}\gnumbox[r]{2}
	&\multicolumn{1}{p{\gnumericColE}|}%
	{\gnumericPB{\raggedleft}\gnumbox[r]{13000}}
\\
	 \multicolumn{1}{|p{\gnumericColB}}%
	{\gnumericPB{\raggedright}\gnumbox[l]{bookshelf}}
	&\gnumericPB{\raggedleft}\gnumbox[r]{19}
	&\gnumericPB{\raggedleft}\gnumbox[r]{13}
	&\multicolumn{1}{p{\gnumericColE}|}%
	{\gnumericPB{\raggedleft}\gnumbox[r]{247}}
\\
	 \multicolumn{1}{|p{\gnumericColB}}%
	{\gnumericPB{\raggedright}\gnumbox[l]{world map}}
	&\gnumericPB{\raggedleft}\gnumbox[r]{3}
	&\gnumericPB{\raggedleft}\gnumbox[r]{1}
	&\multicolumn{1}{p{\gnumericColE}|}%
	{\gnumericPB{\raggedleft}\gnumbox[r]{3}}
\\
	 \multicolumn{1}{|p{\gnumericColB}}%
	{}
	&
	&
	&\multicolumn{1}{p{\gnumericColE}|}%
	{}
\\
\hhline{|----|}
	 \multicolumn{1}{|p{\gnumericColB}}%
	{\gnumericPB{\raggedright}\gnumbox[l]{\textit{Total sum for everything:}}}
	&
	&
	&\multicolumn{1}{p{\gnumericColE}|}%
	{\gnumericPB{\raggedleft}\gnumbox[r]{\textit{1052712}}}
\\
\hhline{|----|}
\end{longtable}

\ifthenelse{\isundefined{\languageshorthands}}{}{\languageshorthands{\languagename}}
\gnumericTableEnd


XFig drawings:

\input{xfigtest.pstex_t}

and PDF pages:

\includepdf[pages=-,angle=22,origin=Bl,width=5cm,height=40mm,keepaspectratio]{../../../lib/examples/beamer-icsi-logo}

If you prefer verbatim input, you can choose
between~\verbatiminput{foo} or~\verbatiminput*{foo}.


\section{URLs and Hyperlinks\index{URL}\nomenclature[www]{URL}{uniform resource locator}}

An URL: \url{http://www.lyx.org} \nomenclature{URL2}{uniform resource locator}

An URL with strange characters: \url{http://www.lyx.org/percent%dollar$}

link:\href{http://www.test.test}{www.test.test}

link2:\href{http://www.test.test}{http://www.test.test}

mail:\href{mailto:www.test.test}{name}

file:\href{file:www.test.test}{www.test.test}

ftp:\href{ftp://www.test.test}{ftp://www.test.test}

ftp2:\href{ftp://www.test.test}{www.test.test}

parser test (stupid, but valid):\href{http://www.test.test}{\}}

parser test (escaped):\href{http://www.test.test}{a brace \} and another one \{ and something}


\section{Lists\index{Lists}}

\listoffigures
\listoftables


\section{Listings}

Inline: \lstinline[language={C++},keywordstyle={\color{green}}]!int a=5;!\\
Float:

\begin{lstlisting}[caption={Example Listing float},label={lst:Example-Listing},language=Python]
# Example listing float
def func(param):
    'this is a python function'
    pass
\end{lstlisting}

Here is an example listing with left line numbering, step ``3'',
language ``Python'', options ``Extended character table'' and
``Space~as~symbol'', range lines 3\,-\,8:

\begin{lstlisting}[extendedchars=true,firstline=3,language=Python,lastline=8,numbers=left,showspaces=true,stepnumber=3]
def func(param):
    'this is a python function'
    pass
def func(param):
'This is a German word: Tsch��'
pass
def func(param):
'this is a python function'
pass
\end{lstlisting}

Special cases:
\begin{lstlisting}[abovecaptionskip=2em,basicstyle={\large\ttfamily},breaklines=true,extendedchars=true,firstline=2,float=h,language={[R/3 3.1]ABAP},lastline=5,numbers=left,numberstyle={\scriptsize},showspaces=true,showstringspaces=false,stepnumber=3,tabsize=4]
hello
\end{lstlisting}

\lstinline[language=TeX]!\begin{centering} hello!


\section{Graphics\index{Graphics}}

There is also some basic support for graphics, in the form
\includegraphics{foo.eps}, or the slightly more elaborate
\includegraphics[bb=10bp 0bp 96bp 96bp,clip,height=1cm, width=1cm]{foo.eps}.

\section{Tables\index{Tables}}

The following example is stolen from the longtable documentation.
Since tex2lyx does not understand the special verbatim code that
was used in the original some lines have been rewritten using
\textbackslash textbackslash etc.
\let\package\textsf
\let\env\textsf
\let\code\texttt
\providecommand\finalclearpage{\clearpage}

\begin{longtable}{@{*}r||p{1in}@{*}}
KILLED & LINE!!!! \kill
\caption
[An optional table caption (used in the list of tables)]
{A long table\label{long}}\\
\hline\hline
\multicolumn{2}{@{*}c@{*}}%
     {This part appears at the top of the table}\\
\textsc{First}&\textsc{Second}\\
\hline\hline
\endfirsthead
\caption[]{(continued)}\\
\hline\hline
\multicolumn{2}{@{*}c@{*}}%
      {This part appears at the top of every other page}\\
\textbf{First}&\textbf{Second}\\
\hline\hline
\endhead
\caption*
{standard foot}\\
\hline
This goes at the&bottom.\\
\hline
\endfoot
\caption*{(last foot)}\\
\hline
These lines will&appear\\
in place of the & usual foot\\
at the end& of the table\\
\hline
\endlastfoot
\env{longtable}  columns  are specified& in the \\
same way as  in the \env{tabular}& environment.\\
\code{@\{*\}r||p\{1in\}@\{*\}}& in this case.\\
Each row ends with a& \code{\textbackslash\textbackslash} command.\\
The \code{\textbackslash\textbackslash} command  has an& optional\\
argument, just as in& the\\
 \env{tabular}&environment.\\[10pt]
See the effect of \code{\textbackslash\textbackslash[10pt]}&?\\
Lots of lines& like this.\\
Lots of lines& like this.\\
Lots of lines& like this.\\
Lots of lines& like this.\\
Also \code{\textbackslash hline} may be used,&  as in \env{tabular}.\\
\hline
That  was a \code{\textbackslash hline}&.\\
\hline\hline
That  was \code{\textbackslash hline\textbackslash hline}&.\\
\multicolumn{2}{||c||}%
{This is a \code{\textbackslash multicolumn\{2\}\{||c||\}}}\\
If a  page break  occurs at a \code{\textbackslash hline} then& a line is drawn\\
at the bottom of one  page  and at the& top of the next.\\
\hline
The \code{[t] [b] [c]} argument of \env{tabular}& can  not be used.\\
The optional argument may be  one of& \code{[l] [r] [c]}\\
to specify whether  the  table  should be& adjusted\\
to the  left, right& or centrally.\\
\hline\hline
Lots of lines& like this.\\
Lots of lines& like this.\\
Lots of lines& like this.\\
Lots of lines& like this.\\
Lots of lines& like this.\\
Lots of lines& like this.\\
Lots of lines& like this.\\
Lots of lines& like this.\\
Lots of lines& like this.\\
Lots of lines& like this.\\
Lots of lines& like this.\\
Lots of lines& like this.\\
Lots of lines& like this.\\
Lots of lines& like this.\\
Lots of lines& like this.\\
Lots of lines& like this.\\
Lots of lines& like this.\\
Lots of lines& like this.\\
Lots of lines& like this.\\
Lots of lines& like this.\\
Some lines may take up a lot of space, like this: &
    \raggedleft This last column is a ``p'' column so this
    ``row'' of the table can take up several lines. Note however that
    \TeX\ will  never break a page within such a row. Page breaks only
    occur between rows of the table or at \code{\textbackslash hline} commands.
    \tabularnewline
Lots of lines& like this.\\
Lots of lines& like this.\\
Lots of lines& like this.\\
Lots of lines& like this.\\
Lots of lines& like this.\\
Lots of lines& like this.\\
Lots of lines& like this.\\
\hline
Lots\footnote{This is a footnote.} of lines& like this.\\
Lots   of   lines& like this\footnote{\env{longtable} takes special
    precautions, so that footnotes may also be used in `p' columns.}\\
\hline
Lots of lines& like this.\\
Lots of lines& like this.
\end{longtable}

From bug 7412 another example with more captions (can currently not produced in LyX):
\begin{longtable}{|l|l|}
\caption{A long table}
\endfirsthead
\caption{A long table -- continued}
\endhead
\multicolumn{2}{r}{{Continued on next page}}
\tabularnewline
\endfoot
\endlastfoot
\hline
\multicolumn{1}{|c|}{Something} & \multicolumn{1}{c|}{Description}\tabularnewline
\hline Lots of lines& like this.\\
\hline Lots of lines& like this.\\
\hline Lots of lines& like this.\\
\hline Lots of lines& like this.\\
\hline Lots of lines& like this.\\
\hline Lots of lines& like this.\\
\hline Lots of lines& like this.\\
\hline Lots of lines& like this.\\
\hline Lots of lines& like this.\\
\hline Lots of lines& like this.\\
\hline Lots of lines& like this.\\
\hline Lots of lines& like this.\\
\hline Lots of lines& like this.\\
\hline Lots of lines& like this.\\
\hline Lots of lines& like this.\\
\hline Lots of lines& like this.\\
\hline Lots of lines& like this.\\
\hline Lots of lines& like this.\\
\hline Lots of lines& like this.\\
\hline Lots of lines& like this.\\
\hline Lots of lines& like this.\\
\hline Lots of lines& like this.\\
\hline Lots of lines& like this.\\
\hline Lots of lines& like this.\\
\hline Lots of lines& like this.\\
\hline Lots of lines& like this.\\
\hline Lots of lines& like this.\\
\hline Lots of lines& like this.\\
\hline Lots of lines& like this.\\
\hline Lots of lines& like this.\\
\hline Lots of lines& like this.\\
\hline Lots of lines& like this.\\
\hline Lots of lines& like this.\\
\hline Lots of lines& like this.\\
\hline Lots of lines& like this.\\
\hline Lots of lines& like this.\\
\hline Lots of lines& like this.\\
\hline Lots of lines& like this.\\
\hline Lots of lines& like this.\\
\hline Lots of lines& like this.\\
\hline Lots of lines& like this.\\
\hline
\end{longtable}

A table*:

\begin{tabular*} % some comment
{0.8\columnwidth}[b]{lr}
two\\
\begin{turn}{-50}lonely\end{turn}&\begin{sideways}lines\end{sideways}
\end{tabular*}

A booktabs table:

\begin{table}[h]
\caption{\label{tab:Special-booktabs-table}Special booktabs-table}


\centering{}%
\begin{tabular}{cccc}
\toprule 
System & Chip\,1 & \multicolumn{2}{c}{Chip\,2}\tabularnewline
\cmidrule(r){2-2}\cmidrule(l){3-4}\morecmidrules \cmidrule{2-4}Detector
thickness in \textmu{}m & 300 & 300 & 700\tabularnewline
\midrule 
Edge angle in \textdegree{} & 3.55 & 2.71 & 7.99\tabularnewline
\addlinespace
Spatial resolution in \textmu{}m & 4.26 & 10.17 & 10.56\tabularnewline
\addlinespace
MTF at $f_{\mathrm{max}}$ & 0.53 & 0.37 & 0.39\tabularnewline
\midrule 
\morecmidrules \cmidrule{3-4}LSF-spatial resolution &  &  & \tabularnewline
in \textmu{}m & 129.7 & 52.75 & 50.78\tabularnewline
in \% of pixel size & 76.3 & 95.9 & 92.3\tabularnewline
\bottomrule
\end{tabular}
\end{table}

A table using dcolumn:

\begin{table} 
\begin{tabular}{ccddd} 
One&Two& 
\multicolumn{1}{c}{\textrm{Three}}& 
\multicolumn{1}{c}{\textrm{Four}}& 
\multicolumn{1}{c}{\textrm{Five}}\\ 
\hline 
one&two&\mbox{three}&\mbox{four}&\mbox{five}\\ 
He&2& 2.77234 & 45672. & 0.69 \\ 
C\footnote{Some tables require footnotes.} 
&C\footnote{Some tables need more than one footnote.} 
& 12537.64 & 37.66345 & 86.37 \\ 
\end{tabular} 
\end{table}


A table using decimal alignment:

\begin{table} 
\begin{tabular}{ccr@{\extracolsep{0pt}.}lr@{\extracolsep{0pt}.}lr@{\extracolsep{0pt}.}l}
One & Two & \multicolumn{2}{c}{Three} & \multicolumn{2}{c}{Four} & \multicolumn{2}{c}{Five}\tabularnewline
\hline 
one & two & \multicolumn{2}{c}{three} & \multicolumn{2}{c}{four} & \multicolumn{2}{c}{five}\tabularnewline
He & 2 & 2&77234  & 45672&  & 0&69 \tabularnewline
C & C & 12537&64  & 37&66345  & 86&37 \tabularnewline
\end{tabular}
\end{table}


\section{Macros}

LyX supports several kinds of macros:
def \def\macroa#1{a #1 a}
global def \global\def\macrob#1{b #1 b}
long def \long\def\macroc#1{c #1 c}
global long def \global\long\def\macrod#1{d #1 d}
providecommand \providecommand{\macroe}[1]{e #1 e}
providecommand* \providecommand*{\macrof}[1]{f #1 f}
newcommand \newcommand{\macrog}[1]{g #1 g}
renewcommand \renewcommand{\macrog}[1]{h #1 h}
newcommand* \newcommand*{\macroi}[1]{i #1 i}
renewcommand* \renewcommand*{\macroi}[1]{j #1 j}
providecommandx \providecommandx{\macrok}[1]{k #1 k}
providecommandx* \providecommandx*{\macrok}[1]{l #1 l}
newcommandx \newcommandx{\macrom}[1]{m #1 m}
renewcommandx \renewcommandx{\macrom}[1]{n #1 n}
newcommandx* \newcommandx*{\macroo}[1]{o #1 o}
renewcommandx* \renewcommandx*{\macroo}[1]{p #1 p}
DeclareRobustCommand \DeclareRobustCommand{\macroq}[1]{q #1 q}
DeclareRobustCommand* \DeclareRobustCommand*{\macror}[1]{r #1 r}
DeclareRobustCommandx \DeclareRobustCommandx{\macros}[1]{s #1 s}
DeclareRobustCommandx* \DeclareRobustCommandx*{\macrot}[1]{t #1 t}

Now use them all:
\macroa{x} \macrob{x} \macroc{x} \macrod{x} \macroe{x} \macrof{x} \macrog{x}
\macroi{x} \macrok{x} \macrom{x} \macroo{x} \macroq{x} \macror{x}
The following tow don't work, but they should???
%\macros{x} \macrot{x}


\section{Special formattings\index{Special formattings}}

\subsection{LyX line}

test

{\tiny \lyxline{\tiny}}{\tiny \par}

test {\Huge Test} {\tiny test} test

\lyxline{\normalsize}

test

{\Huge \lyxline{\Huge}}{\Huge \par}


\subsection{Rules}

test

{\tiny \rule[0.5ex]{1\linewidth}{1pt}}{\tiny \par}

test\rule[0.001\textwidth]{1\columnwidth}{0.05\textheight}

test

\rule[0.5ex]{1\columnwidth}{1pt}

\noindent \rule{1ex}{0.5cc}test

\textcolor{red}{\rule[-4ex]{5in}{1cm}}


\subsection{Phantoms}

test\phantom{Wow}test

test\hphantom{\textcolor{green}{W}\textbf{ow}\textsuperscript{Wow}}test

test\vphantom{Wow\textbackslash{}\&\%}test

$test\hphantom{Wow}test$

\[
test\phantom{\textcolor{green}{W}\mathbf{ow}\textsuperscript{Wow}}test
\]


$test\vphantom{Wow\&\%}test$


\subsection{TIPA}

\begin{IPA}
\textsubring{e}
\r{e}
\textsubwedge{e}
\textsubumlaut{e}
\textsubtilde{e}
\textseagull{e}
\textsubbridge{e}
\textinvsubbridge{e}
\textsubsquare{e}
e\textsuperscript{h}
\textsubrhalfring{e}
\textsublhalfring{e}
\textsubplus{e}
\=*{e}
(don't mix up with \b{e})
\"{e}
\textovercross{e}
\s{e}
\textsyllabic{e}
\textsubarch{e}
e\textrhoticity
e\textsuperscript{w}
e\textsuperscript{j}
e\textsuperscript{\textgamma}
e\textsuperscript{\textrevglotstop}
\textsuperimposetilde{e}
\textraising{e}
\textlowering{e}
\textadvancing{e}
\textretracting{e}
\~{e}
e\textsuperscript{n}
e\textsuperscript{l}
e\textcorner{}
\H{e}

\`{e}
\={e}
\'{e}
\textdoublegrave{e}
\v{e}
\^{e}
\texthighrise{e}
\textlowrise{e}
\textrisefall{e}
\u{e}
\end{IPA}

\textipa{pbtd\:t\:dc\textbardotlessj k\textscriptg q\;G}
\textipa{PmMn\:n\textltailn N\;N\;Br\;RR}
\textipa{\:rFBfvTDszSZ\:s}
\textipa{\:z�JxGXK\textcrh QhH\textbeltl{}}
\textipa{\textlyoghlig V\*r\:Rh\textturnmrleg l\:lL\;L}

\textipa{\!o|!\textdoublebarpipe ||\!b\!d\!j\!g\!Ge'}

\textipa{iy1\textbaru W\textcolor{red}{uIY}Ue�987o@E�3\textcloserevepsilon 2O�5a\textscoelig A6}

\textipa{\*ww4\;H\textbarrevglotstop\textbarglotstop C\textctz\textturnlonglegr\texththeng \texttoptiebar{ar}\textbottomtiebar{tz}}

\textipa{:;e\textprimstress\textsecstress \textvertline{}\textdoublevertline{}.\t*{ }}

\textipa{\tone{55}\tone{44}\tone{33}\tone{22}\tone{11}\textdownstep\textupstep \tone{15}\tone{51}\tone{45}\tone{12}\tone{454}\textglobrise \textglobfall{}}

TIPA code in math:
$\textipa{\tone{55}|\text{\!b{}\!d{}\!g{}\!G{}\textglobfall{}\textvertline{}\textdoublevertline{}}}$

\subsection{Other symbols}
All three should be converted to U+00C5: \AA{} \r A \r{A} (not U+212B).
All three should be converted to U+00E5: \aa{} \r a \r{a}.

\subsection{Line breaks\sindex[breaks]{Line breaks}}

They can also or be broken by a newline\\
or by a starred newline \\*
or by a newline with space, comment and argument \\ %hu
[3cm]
or by a newline command \newline
or by a line break \linebreak
or by a defined line break \linebreak % again with a comment
	[4]

There are even newlines with weird arguments, but these are not
handled by LyX\\*[1cm]
so we try to use ERT in this case.

\subsection{Page breaks\sindex[breaks]{Page breaks}}

They can also or be broken by a newpage \newpage
or by a page break \pagebreak
or by a defined page break \pagebreak % again with a comment
	[4]


\section{Special characters\index{Special characters}}

Then one has those macros with a long name for a short meaning, like
\textasciitilde, \textasciicircum{} or \textbackslash{}, \slash{},
\nobreakdash- and the characters
that LaTeX wants to espace because they are active, like \_\&\#\$\{\}\%.

And what about special characters like hyphe\-nation mark,
ellipsis\ldots, and end-of-sentence\@. LyX also supports a menu
separator\lyxarrow{}and a spif\textcompwordmark{}fy ligature break.

There are dashes: endash in short form -- and long form \textendash,
emdash is alike: --- and \textemdash. If we really want several hyphens
in a row, we need to separate them: -{}-, -{}-{}-, -{}-{}-{}- etc.

LyX translates the phrases LyX, TeX, LaTeX2e and LaTeX
to the commands \LyX{}, \TeX{}, \LaTeXe{} and \LaTeX{}.
If these phrases occur as part of other words (like 1LyX or aTeX or LaTeX3)
they should not be put into ERT.

Test for whitespace handling of commands: The following lines should
result in identical output:

builtin \textasciicircum{} unicodesymbols \j{} user \mycommand{} xx\par
builtin \textasciicircum {} unicodesymbols \j {} user \mycommand{} xx\par
builtin \textasciicircum % with a comment
{} unicodesymbols \j % and a second one
{} user \mycommand % and another
{} xx

A sub\textsubscript{sc\emph{ript}} and super\textsuperscript{script
with $a^2+b^2=c^2$ math}.

latex2writer outputs quotes in braces like {\textquoteright}. We swallow the,
but this was one broken (bug 8903).


\section{Mathematics\index{Mathematics}}

Let $f:\left[  a,b\right]  \rightarrow%
%TCIMACRO{\U{211d} }%
%BeginExpansion
\mathbb{R}
%EndExpansion
$.

\[
\begin{array}{rclccc}
1 + 2 & = & 3 
\multicolumn{3}{c}{4 < 5 \leq 6 }
\end{array}
\]


\section{Lists/Indices}

single \textbackslash{}addcontentsline: \addcontentsline{toc}{section}{test}

with \textbackslash{}addcontentsline and \textbackslash{}phantomsection:

\bibliographystyle{test}
\phantomsection\addcontentsline{toc}{section}{\refname}\bibliography{xampl}


normal:

\bibliographystyle{test}
\bibliography{xampl}


with wrong \textbackslash{}addcontentsline:

\bibliographystyle{test} \addcontentsline{toc}{section}{test} \bibliographystyle{test}
\bibliography{xampl}


with \textbackslash{}addcontentsline and \textbackslash{}nocite\{{*}\}:

\bibliographystyle{test}
\addcontentsline{toc}{section}{\refname}\nocite{*}
\bibliography{xampl}


undefined index:

\printindex{}

index "idx":

\printindex[idx]{}

index "new":

\printindex[new]{}

subindex "new":

\printsubindex[idx]{}

index of all indices:

\printindex*{}

subindex of all indices:

\printsubindex*{}

normal nomenclature:

\printnomenclature hello

manually set width: \settowidth{\nomlabelwidth}{URL2} \printnomenclature{}

nomenclature with set width:

\printnomenclature[0.02\linewidth]{}

\end{document}
